\documentclass[11pt]{article}

\usepackage[margin=1in]{geometry}
\usepackage{amsmath,amssymb,amsthm,mathtools}
\usepackage{mathrsfs}
\usepackage{bm}
\usepackage{enumitem}
\usepackage{booktabs}
\usepackage[dvipsnames]{xcolor}
\usepackage[most]{tcolorbox}
\usepackage{microtype}
\usepackage[hidelinks]{hyperref}
\usepackage[nameinlink]{cleveref}

\allowdisplaybreaks
\setlength{\parskip}{0.4em}

%==================== macros ====================
\newcommand{\E}{\mathbb{E}}
\newcommand{\C}{\mathbb{C}}
\newcommand{\Xcal}{\mathcal{X}}
\newcommand{\Lcal}{\mathcal{L}}
\newcommand{\Ccal}{\mathscr{C}}
\newcommand{\Om}{\Omega}
\newcommand{\Qcoll}{\mathfrak{Q}_{\mathrm{Coll}}}
\DeclareMathOperator{\tr}{tr}
\DeclareMathOperator{\hol}{hol}
\newcommand{\epi}[1]{{\small\sffamily\textbf{[#1]}}}

% headline / key-statement box (prose)
\newtcolorbox{keystatement}{%
  enhanced, sharp corners, boxrule=0.6pt,
  colback=black!4, colframe=black!70,
  left=8pt, right=8pt, top=6pt, bottom=6pt,
  before skip=8pt, after skip=8pt}

% titled box for the central conjectural principle
\newtcolorbox{principlebox}[1]{%
  enhanced, sharp corners, boxrule=0.8pt,
  colback=black!3, colframe=black!75,
  coltitle=black, fonttitle=\bfseries,
  title={#1},
  left=8pt, right=8pt, top=6pt, bottom=6pt,
  before skip=10pt, after skip=10pt}

\theoremstyle{remark}
\newtheorem*{remark}{Remark}

%==================== title ====================
\title{The Quenched Residue--Height Spectral Shadow}
\author{}
\date{June 9, 2026}

\begin{document}
\maketitle

\begin{remark}[Status of these notes]
These are working research notes. The construction below separates objects that are
already finite and checkable from the one statement that remains a new, unproven theorem;
that separation is the whole point of the note, and it is preserved as written. The central
\emph{Quenched Spectral-Shadow Principle} is conjectural (the text treats it as something
that ``would'' be proved). Project-internal terms (notably \textsf{H23a}, the
quasi-stationary distribution (QSD) machinery, the residue spaces $\Om_Q$, and the exact
affine correction $\varepsilon_{\mathrm{exact}}$) are used without restatement because they
refer to a larger body of work. Nothing here has been verified.
\end{remark}

\section{Keeping the pieces that survive}

The useful pieces of the preceding answers are not the proposed names themselves:
``Gibbs measure,'' ``arithmetic coboundary sheaf,'' ``spectral-coboundary module,'' or
``Collatz Selberg zeta function.'' Those are fingers pointing toward the missing object,
not the object itself. The pieces worth keeping are the following.

\section{The annealed thermodynamic formalism}

First, the annealed theory is genuinely unified by the transfer-operator / thermodynamic
formalism of the Collatz height cocycle. The basic height increment is
\[
\Delta_b = \log_2 3 - b, \qquad b = v_2(3x+1),
\]
and under the Haar/annealed valuation law
\[
p_b = 2^{-b}, \qquad b\ge 1,
\]
the moment transform is
\[
M(s) = \sum_{b\ge1} 2^{-b}\, 2^{s(\log_2 3 - b)} = \frac{3^s}{2^{1+s}-1}.
\]
The Cram\'er root is
\[
M(1) = 1,
\]
and the critical tilted law is
\[
p_b^{(1)} = 2^{-b}\, 2^{\log_2 3 - b} = 3\cdot 4^{-b} = \tfrac34\left(\tfrac14\right)^{b-1}.
\]
This package explains the annealed height tail, the Cram\'er tilt, the conditioned
high-ascent law, and the pressure side of the killed-walk/QSD machinery.

\section{Finite residue-height transfer operators}

Second, the finite residue-height transfer operators are the safe spectral objects. For a
finite modulus $Q$ and strip width $L$, one can define a finite state space
\[
\Xcal_{Q,L} = \Om_Q \times I_L,
\]
where $\Om_Q$ is the appropriate residue space and $I_L$ is the discretized height/deficit
strip. Edges are exact admissible inverse-Syracuse transitions, indexed by valuations $b$,
with residue update
\[
a' \equiv 3^{-1}\bigl(2^b a - 1\bigr) \pmod Q,
\]
and height update
\[
z' = z + \log_2 3 - b + \varepsilon_{\mathrm{exact}}(a,z,b),
\]
where $\varepsilon_{\mathrm{exact}}$ records the exact affine correction when it is
retained. Killing occurs when $z'$ exits the strip.

For each $s$, the finite tilted transfer operator is
\[
(\Lcal_{s,Q,L} f)(x) = \sum_{e\,:\,y\to x} w_s(e)\, f(y),
\]
with edge weight of the form
\[
w_s(e) = m(e)\, 2^{s(\log_2 3 - b(e))}.
\]
Here $m(e)$ is the exact finite-fiber weight; in the pure annealed model it is
$2^{-b(e)}$. The finite determinant
\[
\Xi_{Q,L}(z,s) = \det\!\bigl(1 - z\Lcal_{s,Q,L}\bigr)
\]
is unambiguously defined. Its logarithm records closed paths in the finite residue-height
graph:
\[
\log \Xi_{Q,L}(z,s)^{-1} = \sum_{n\ge1} \frac{z^n}{n}\, \tr\!\bigl(\Lcal_{s,Q,L}^{\,n}\bigr).
\]
This is the safe finite analogue of a Ruelle or Selberg zeta function. A global
$S$-arithmetic Selberg zeta function is not presently known and should not be assumed.

\section{Character sectors and the coboundary obstruction}

Third, finite character sectors are the natural home of the coboundary obstruction.
Decompose the finite operator by residue characters:
\[
\Lcal_{s,Q,L} = \Lcal_{s,Q,L,0} \oplus \bigoplus_{\chi\ne 1} \Lcal_{s,Q,L,\chi}.
\]
A nontrivial character $\chi$ is a rank-one local system on the finite residue graph.
\textsf{H23a}-type statements assert a uniform nontrivial-character gap:
\[
\rho\bigl(\Lcal_{s,Q,L,\chi}\bigr) \le (1-\kappa_Q)\,\rho\bigl(\Lcal_{s,Q,L,0}\bigr),
\qquad \chi\ne 1.
\]
Equivalently, the finite residue graph has no nontrivial local system whose holonomy is
invisible to the dynamics. This is the rigorous finite shadow of the phrase
``spectral-coboundary obstruction.''

\section{The quenched residue-height spectral shadow}

The object suggested by these pieces is therefore not a single Gibbs measure, not a
completed sheaf, and not an already-existing zeta function. It is the following pro-object:
\[
\boxed{\;
\Qcoll = \Bigl\{\, \Xcal_{Q,L},\ \Lcal_{s,Q,L},\ \Lcal_{s,Q,L,\chi},\
\Xi_{Q,L,\chi},\ \Ccal_{Q,L} \,\Bigr\}_{Q,L}
\;}
\]
Here $\Ccal_{Q,L}$ denotes the finite local-system/coboundary package generated by the
nontrivial residue characters and their holonomies on the exact residue-height graph. More
explicitly,
\[
\Ccal_{Q,L} = \bigoplus_{\chi\ne 1} \C_\chi,
\]
viewed as a coefficient system over the finite directed graph $\Xcal_{Q,L}$. A closed path
$\gamma$ has holonomy
\[
\hol_\chi(\gamma) = \prod_{e\in\gamma} \chi(e),
\]
and the finite certificate says that for every nontrivial $\chi$, some certified cycle
satisfies
\[
\hol_\chi(\gamma) \ne 1.
\]
Thus nontrivial finite characters cannot behave as coboundaries on the primitive bulk
component.

We call the pro-object $\Qcoll$ the \emph{quenched residue-height spectral shadow}. It is a
shadow because it is not yet the missing theorem itself. It is the finite system in which
any future quenched obstruction must become visible.

\section{The Quenched Spectral-Shadow Principle}

The desired inverse principle can now be stated cleanly.

\begin{principlebox}{Quenched Spectral-Shadow Principle (conjectural)}
If a deterministic positive-integer orbit violates the annealed height/valuation law
persistently, for example by producing an infinite sub-equilibrium tail, repeated critical
high ascents, or sustained deviation of
\[
D_m - m\log_2 3
\]
from the equilibrium line, then there exists a finite pair $(Q,L)$ and a nontrivial
character sector $\chi\ne 1$ such that the bad orbit has a nonzero finite spectral shadow in
$\Ccal_{Q,L}$. Equivalently, the bad orbit forces an approximate peripheral obstruction:
\[
\rho_{\mathrm{eff}}\bigl(\Lcal_{s,Q,L,\chi}\bigr)
\approx
\rho\bigl(\Lcal_{s,Q,L,0}\bigr).
\]
\end{principlebox}

If this principle were proved, \textsf{H23a} would rule out the obstruction:
\[
\rho\bigl(\Lcal_{s,Q,L,\chi}\bigr) \le (1-\kappa_Q)\,\rho\bigl(\Lcal_{s,Q,L,0}\bigr),
\qquad \chi\ne 1.
\]
Thus the contradiction would be:
\[
\text{bad deterministic orbit}
\;\Longrightarrow\;
\text{finite spectral/coboundary shadow}
\;\Longrightarrow\;
\text{violation of \textsf{H23a}}.
\]
This is the precise form of the missing annealed-to-quenched bridge.

\section{Finite versus new}

The point of the construction is that it separates what is already finite and checkable
from what remains genuinely new.

\begin{center}
\renewcommand{\arraystretch}{1.3}
\begin{tabular}{@{}p{0.40\linewidth} p{0.52\linewidth}@{}}
\toprule
\textbf{Component} & \textbf{Status}\\
\midrule
$\Lcal_s,\ M(s),\ p_b^{(s)}$ & annealed thermodynamic object\\
$\Lcal_{s,Q,L}$ & finite exact residue-height operator\\
$\Xi_{Q,L}(z,s)$ & finite determinant / zeta shadow\\
$\Ccal_{Q,L}$ & finite character / coboundary local-system package\\
\textsf{H23a} gap & finite spectral obstruction exclusion\\
Quenched Spectral-Shadow Principle & new missing theorem\\
\bottomrule
\end{tabular}
\end{center}

\section{Discarded overclaims}

The discarded overclaims are equally important. A Gibbs or SRB measure does not solve the
problem, because ergodicity gives almost-everywhere statements, not control of every
arithmetically prescribed integer orbit. A global $S$-arithmetic Selberg zeta function is
not presently constructed, and no known trace formula applies to the noninvertible Collatz
monoid. A sheaf or coboundary module is only meaningful after a finite residue-height graph,
local systems, and restriction maps are specified. Baker theory may eliminate certain
cycle-like Diophantine resonances, but it does not by itself prove deterministic
equidistribution. Exceptional Lie theory currently supplies no load-bearing object.

\section{The object we can construct now}

Thus the answer to the question ``what is the object?'' is:

\begin{keystatement}
The object we can construct now is the pro-system of finite residue-height tilted transfer
operators, their character-sector decompositions, their finite Fredholm determinants, and
their local-system/coboundary obstruction package:
\[
\Qcoll = \bigl\{\,(\Xcal_{Q,L},\ \Lcal_{s,Q,L},\ \Ccal_{Q,L})\,\bigr\}_{Q,L}.
\]
It is the finite spectral shadow in which any bad deterministic orbit must be detected if
the annealed-to-quenched inverse theorem is true.
\end{keystatement}

This is the finger, not the moon. The moon would be the theorem
\[
\text{bad quenched orbit}
\;\Longrightarrow\;
\text{nontrivial finite spectral shadow}.
\]
The finger tells us where to look: not at a free-standing Gibbs measure, not at an
unconstrained symbolic shift, and not at speculative global zeta functions, but at the
finite residue-height spectral shadows of deterministic orbits and the missing inverse
theorem that would force every bad orbit to cast one.

\end{document}
